Data Characterisitcs
The data used has been cleaned and rebased at approriate time period to allow for a better comparison between the nations, removing the influence of other factors such as currency exchange.
Energy price cpi - the average price of energy given in monthly increments, rebased at 2015=100 to aid in comparision between the three nations.

Headline cpi - the monthly all-items, non-adjusted consumer price index, rebased at 2002=100 to aid in comparision between the three nations.

Domestic consumption - the total domestic consumption of energy, given in quarterly increments, rebased at 2015=100 to aid in comparision between the three nations.

Energy supply - the type of energy production method used within the constituent nations.

Scatter Plots
CPI Energy Price vs Headline CPI
This scatter plot shows the relationship between the energy price cpi and the headline cpi for the three nations. Each point represents a monthly observation, with the x-axis representing the energy price cpi and the y-axis representing the headline cpi. 
The plot indicates a positive correlation between the two variables, suggesting that as energy prices increase, headline inflation also tends to increase.
General positve direction, although the year 2022 has some variation.

CPI Energy vs CPI Headline Slope
This plot is an amalgamtion of the previous scatter plots for all the nations with the addiotn of the slope, which shows the corresponding change in headline CPI for every one unit rise in energy price.
The United States and Canada have similar slopes, whcih could roughly be expected due to their geographic location and economic strucutre.
The United Kingdom has a lower slope indicating that headline CPI absorbed less of the energy shcok per unit.
High R^2 values for each of the nations which needs to be further interogated. 

CPI Headline vs Con
This scatter plot illustrates the relationship between the headline cpi and domestic consumption. Each point represents a quarterly observation, with the x-axis representing the energy price cpi and the y-axis representing domestic consumption.
The slope indicates the change in domestic consumption per one unit change in headline CPI.
Canada had the highest slope which was unexpected, meaning that there was strong growth in domestic consumption alongside the growth in headline CPI.
The United states had a nearly 1:1 relationship between the two variables, which was simialry unexpected.
The United Kingdom faired worese than both the US and Canada, with a much lower slope, although this was not unexpected.

CPI Headline vs Con EQUAL
Gives the same scatter plots as CPI Headline vs Con, but with only the data for years in which all three nations have data available for.
In the equalised version Canada and the United States have similar slopes near the 1:1 ratio, whislt the United Kingdom's remained relatively low at 0.664, shwoing that there was a meaninful difference in the growth between the two variabels.
This time period is more recent and shows changes between the non-equlaised version, it may be useful to include economic growth rates within the model to try and control for that factor.

Correlation Maps
CPI Energy vs CPI Headline Correlation
Shows the correlation between the energy price cpi and the headline cpi for each of the three nations. The correlation coefficient is calculated to quantify the strength and direction of the relationship between the two variables. A positive correlation coefficient indicates that as energy prices increase, headline inflation also tends to increase.
All three nations have a correaltion close to one in the year 2022, showing that during that year there is an almost perfect correltion between the headline CPI and Energy CPI.
The United Kingdom and Canada have a higher median correlation value than the United States inidcating that they're generally more affected by energy price swings transmitting into headline CPI than the UNited States, which was to be expected.

CPI energy vs CPI Headline LAG
Show how stronlgy a change in the energy CPI at a given time period predicts a change in the headline CPI.
The positve value shown indicate that energy CPI leads the headline CPI, demonsatrating that changes in energy prices can be a predictor of future changes in headline inflation.
These variables have been tested to see if theyre statistically significant, with the results showing that the relationship between energy prices and headline inflation is indeed significant for most of the measured values in Canada and the United Kingdom.
The United States is the only nation to not have many significant variables which warrants further investigation into their policy strategies at the time and how they're economy reacted to the shock.

Energy CPI vs CPI Headline Indivdual
Shows the same as the CPI Energy vs CPI Headline, but with the three countries seperated out instead of within the same graph for easier reference and identification of any signifcant chganges and events present wothin the graph.

Energy Production All
Shows the methods of energy production over tiem between the three nations, showing how these methods changed and the percentage values of the different methods used.
It shows the total amount of energy produced within the economy's and what share of it was produced by which type in a given year, with 2022 specifically shown as its the year of the examined energy crisis.

Correlation Analysis
A first order effect does exist between the energy prices and the headline inflation rate as the correaltion between the two variables is positve and relatively large.
The slope when comparing the two variables is simialry positve and relatively high, shwoing that there is a positve first order effect between the two variabels, where one increase so does the other.
The Simpsons Paradox is certainly worth considering, especailly as the three countries do not share the same systems of governance or perhaps repsonce to the crisis, an immediate example would be the United Kingdom's usgae of an energy price cap. This should be included in the futre as a comparison and to see whether the price cap was reached.
